%
% Title page for Thesis.
%

\vspace*{1cm}
\begin{center}



{\Huge An agent-based modelling approach to spatial economic theory\vspace{1.5cm}\\}

{\Large by\\Dan Olner\vspace{1.5cm}\\




Submitted in accordance with the requirements for the degree of Doctor of Philosophy\vspace{2cm}\\

%\vspace{4cm}\\

\null
\vfill

University of Leeds\\ 
School of Geography\\
March 2014\\
%March 2014\vspace{2cm}\\
}


\end{center}

\clearpage{\pagestyle{empty}\cleardoublepage}

\thispagestyle{empty}

%\begin{flushbottom}
\null
\vfill

The candidate confirms that the work submitted is his own and that appropriate credit has been given where reference has been made to the work of others.

%The candidate confirms that the work submitted is his/her own, except where work which has formed part of jointly-authored publications has been included. The contribution of the candidate and the other authors to this work has been explicitly indicated overleaf. The candidate confirms that appropriate credit has been given within the thesis where reference has been made to the work of others.
\vspace{1cm}\\

%The candidate confirms that the work submitted is his own and that
%appropriate credit has been given where reference has been made to
%the work of others.\vspace{1cm}\\
%\begin{center}

This copy has been supplied on the understanding that it is copyright material and that no quotation from the thesis may be published without proper acknowledgement.

%\end{center}

\vspace{1cm}\\

%The candidate confirms that the work submitted is his own and that
%appropriate credit has been given where reference has been made to
%the work of others.\vspace{1cm}\\
%\begin{center}

The right of Dan Olner to be identified as Author of this work has been asserted by him in accordance with the Copyright, Designs and Patents Act $1988$.

\begin{center}

\vspace{2cm}\\

%This  work  is  licensed  under  the  Creative  Commons Attribution­NonCommercial-ShareAlike $3.0$  Unported  License. 
%To  view  a  copy  of  this  license,  visit http://creativecommons.org/licenses/by-­nc-­sa/3.0/
This  work  is  licensed  under  the  Creative Commons Attribution-NonCommercial-ShareAlike $2.0$ License (United Kingdom). 
To  view  a  copy  of  this  license,  visit http://creativecommons.org/licenses/by-nc-nd/2.0/uk/

%\copyright 2014 The University of Leeds and Dan Olner

\end{center}

%\end{flushbottom}

%\end{center}



%End of title page



\clearpage{\pagestyle{empty}\cleardoublepage}


\thispagestyle{empty}

%\hspace*{-0.75cm} Section \ref{sec:EthnicConcentration} in Chapter \ref{chapter:5} is based on the joint publication:\\\\
%{\bf Harland, K. and Stillwell, J. (2007)} `Evidence of ethnic minority dispersal in Leeds', \\\hspace*{1cm}\emph{The Yorkshire and Humber Regional Review}, 17(2): 21-23\\\\
%This work appears as published in Appendix \ref{appendix:2} for reference.  I declare that the research for this publication was solely my own work and that I am the lead author.  The contribution of the second named author, John Stillwell, was purely editorial and advisory.

%end of included material
\clearpage{\pagestyle{empty}\cleardoublepage}


\thispagestyle{empty}

\vspace*{3cm}

\begin{center}
{\Large For Sue\vspace{1.5cm}}\\


\vspace{1cm}

%\begin{quote}
{\large
\emph{Paper cup of a boat\\
Heaving chest of the sea\\
Carry both of us\\
Carry her, carry me}}\\
\vspace{1cm}
\hspace*{4cm}Elbow}

%\end{quote}
\end{center}


%\clearpage{\pagestyle{empty}\cleardoublepage}


%\thispagestyle{empty}

%\vspace*{5cm}

%\begin{quote}
%{\large
%``Economic ideas are always and intimately a product of their own time and place; they cannot be seen apart from the world they interpret. And that world changes - is, indeed, in a constant process of transformation.''\\
%\hspace*{8cm}\emph{\citealt{galbraith_history_1991}, pp.1-2}}
%}
%\end{quote}

%\vspace{5cm}

%\begin{comment}
%\begin{quote}
%{\large
%``There is nothing new under the sun.''\\% Is there anything of which one can say, `Look! This is something new'?'
%\hspace*{8cm}\emph{Ecclesiastes 1:9,10}}
%}
%\end{quote}
%\end{comment}


\clearpage{\pagestyle{empty}\cleardoublepage}

%Start Acknowledgements


\chapter*{Acknowledgements}
\thispagestyle{empty}

%\par
First, to my ever-patient supervisors, Andrew Evans and Alison Heppenstall: you managed the supervision equivalent of herding cats. Thank you so much for getting me through to the end in one piece. Thank you to Emanuel Gloor and Kirk Harland, who at different times carried out supervision tasks both before and after officially `joining the team'. My research support group managed the unenviable job of making me focus in the right direction: Graham Clarke, Stuart Hodkinson and Evan Fraser. Thanks also to Evan for inviting me to join discussion groups in the School of Earth and Environment at Leeds. Members of the `Multi-agent Simulation Society' in the School of Geography have offered wonderful insights and ideas over the years, and listened to many rambling presentations: thank you. Thanks to Malcolm Sawyer for so willingly offering advice on the economic issues, and to Andrew Neil Ross for putting up with my maths questions. Thank you to Matthew Handy for stoically dealing with my maths problems and for pointing me in the right direction. Andy Goldring's enthusiasm helped me on many occasions to remember why I was doing this, and fed me a plethora of ideas. My office mates made up an unofficial PhD support group: thank you for everything. Thank you to everyone in the School of Geography, and especially to Jacqui Manton who, among other things, made it possible for me to spend time working in Australia. Thank you to Professors Mark Birkin and Seth Bullock for examining the thesis and especially for all your support getting me through the final stages. And to every one of my friends who got me through those last unexpected months: I don't know where I would have been without you. Lastly, thank you to the ESRC for giving me the opportunity to do this PhD: I'm glad I don't have to do it again, but it's been the most amazing experience.

\vspace{1cm}

Lastly, this thesis is dedicated to Susannah Bird. For years, you never failed to be there for me. You were my unofficial supervisor and talked me through seemingly impenetrable PhD thickets. The many, many times when I couldn't see a way forward, you told me to trust you. You were heroically patient and kind, and you worked so hard to provide me with the space to get it done. Every day, you told me: `you \emph{can} do it'. You never stopped believing in me, even when I argued with you and said you shouldn't. I won't forget.

\clearpage{\pagestyle{empty}\cleardoublepage}

%Start Abstract

\chapter*{Abstract}
\thispagestyle{empty}

This thesis presents an agent-based modelling framework, leading to a `transmission belt' model connecting supply and demand in a spatial setting, as a step towards understanding the ``complex, integrated, tightly-coupled global fabric of exchange'' (\citealt{korowicz_tipping_2010} p.2) that makes up the modern spatial economy. Informed by a careful unpacking of economic modelling ideas, the `transmission belt' model succeeds in producing a stable equilibrium of consumption and production across space. The motivation for creating the model framework is as follows. While some argue `it is better to assume that moving goods is essentially costless' (\citealt{glaeser_cities_2004}) because space costs are so low, others point out that a future of `oil depletion' (\citealt{sorrell_assessment_2009}) and high carbon prices due to climate change mean that the implications of future cost changes are not well understood. The thesis examines the most prominent economic approach to space costs, geographical economics (GE), which finds ways to avoid key modelling problems imposed by space but has a tight `mathematical straightjacket' (\citealt{martin_new_1999}) of assumptions. This thesis keeps to simple utility functions for describing actors' preferences, but uses an agent-based modelling approach to break out of these assumptions. While spatial agent models have dealt with a huge range of actor-environment interactions, very few examine traditional spatial economic problems. As a consequence, simple but powerful spatial economic ideas have been neglected. Much can be learned from issues that faced economists throughout the twentieth century. By closely examining these ideas, the thesis asks: what obstacles have stopped agent-based modelling from tackling the `big questions' of spatial economics? The answers suggested are: a lack of research directly tackling the space of `dependencies' in agent model development and a broader sense that it has broken all ties with history and can learn nothing from past modelling efforts.

%End Abstract